\section*{الفصل \ref{ch:b1:electronic-effects} --- التأثيرات الإلكترونية والوسائط التفاعلية}

\begin{solution}{exo:b1:electronic-effects:1}
2-ميثيل البوتان، \ce{CH3-CH(CH3)-CH2-CH3}: مجموعات \ce{CH3} الثلاث أولية،
و\ce{CH2} ثانوية، وCH ثالثية. 2-برومو-2-ميثيل البروبان: مجموعات
\ce{CH3} الثلاث أولية، والكربون المركزي ثالثي. ومن البروميدات الأربعة
\ce{C4H9Br} (1-برومو البوتان، و2-برومو البوتان، و1-برومو-2-ميثيل البروبان،
و2-برومو-2-ميثيل البروبان)، يعطي الأخير \omterm{def:b1:electronic-effects:intermediates}{الكاتيون الكربوني} الثالثي، وهو
الأكثر ثباتًا.
\end{solution}

\begin{solution}{exo:b1:electronic-effects:2}
$10^{4.76 - 2.87} = 10^{1.89} = 78$؛ $10^{4.76 - 0.52} = 10^{4.24} =
\num{1.7e4}$.
\end{solution}

\begin{solution}{exo:b1:electronic-effects:3}
الإيثانوات: تتبادل ذرتا الأكسجين C=O و\ce{C-O^-}.
4-نيتروفينوكسيد: يكوّن \omterm{def:b1:lewis-resonance-vsepr:lewis}{الزوج الحر} للأكسجين رابطة C=O، وتنزاح الروابط الثنائية للحلقة،
ويكوّن الكربون الحامل للمجموعة \ce{NO2} رابطة C=N، ويذهب زوج من N=O
إلى الأكسجين: بنية فيها C=O في أحد طرفي الحلقة وC=N في الطرف الآخر
وذرتا أكسجين النيترو كلتاهما سالبتان (بنية كينونية). الكاتيون الأليلي:
\ce{CH2=CH-CH2+} $\leftrightarrow$ \ce{+CH2-CH=CH2}.
\end{solution}

\begin{solution}{exo:b1:electronic-effects:4}
السهم الثاني: من الرابطة O--H في الماء إلى أكسجينه؛ والنواتج \ce{H2O}
(من \ce{OH-}، وقد صار متعادلًا) و\ce{OH-} (يحتفظ أكسجين الماء السابق
بالزوج، والشحنة $-1$). وفي الأمونيا: سهم من \omterm{def:b1:lewis-resonance-vsepr:lewis}{الزوج الحر} لذرة N إلى
ذرة H من \ce{H3O+}، وآخر من تلك الرابطة O--H إلى الأكسجين: فتصير شحنة N
مساوية $+1$، وشحنة O صفرًا.
\end{solution}

\begin{solution}{exo:b1:electronic-effects:5}
\babelsublr{\foreignlanguage{arabic}{الإيثانول (15.9)} $<$ \foreignlanguage{arabic}{الماء (14.0)} $<$ \foreignlanguage{arabic}{الفينول (9.99، عدم تمركز في الحلقة)} $<$
\foreignlanguage{arabic}{4-نيتروفينول (7.16، تأخذ \ce{NO2} الشحنة)} $<$ \foreignlanguage{arabic}{حمض الإيثانويك (4.76،
ذرتا أكسجين متكافئتان)} $<$ \foreignlanguage{arabic}{حمض ثلاثي كلورو الإيثانويك (0.51، ثلاث ذرات كلور
ذات تأثير $-I$)}}.
\end{solution}

\begin{solution}{exo:b1:electronic-effects:6}
\babelsublr{\foreignlanguage{arabic}{الأنيلين (4.60)} $<$ \foreignlanguage{arabic}{البيريدين (5.23)} $<$ \foreignlanguage{arabic}{الأمونيا (9.25)} $<$ \foreignlanguage{arabic}{الميثيل أمين
(10.66)}}. في الأنيلين يكون \omterm{def:b1:lewis-resonance-vsepr:lewis}{الزوج الحر} مشتركًا مع الحلقة
(بنى فيها C=\babelsublr{N$^+$} والشحنة السالبة في الموضعين \textit{ortho}
و\textit{para})؛ وبرتنة N تكلّف خسارة عدم التمركز هذا.
\end{solution}

\begin{solution}{exo:b1:electronic-effects:7}
في أيون 3-نيتروفينوكسيد، تضع \omterm{def:b1:lewis-resonance-vsepr:resonance}{بنى الرنين} الشحنة على
الكربونين \textit{ortho} و\textit{para} بالنسبة إلى الأكسجين، وأبدًا على الكربون
الحامل للمجموعة \ce{NO2}: فلا عون ميزوميري، ولذلك هو أضعف من المتماكب 4. لكن
مجموعة النيترو تبقى تسحب الإلكترونات عبر الروابط $\sigma$ ($-I$)،
ولذلك هو أقوى من الفينول.
\end{solution}

\begin{solution}{exo:b1:electronic-effects:8}
\babelsublr{\ce{CH3+} $<$ \ce{CH3CH2+} $<$ \ce{CH2=CH-CH2+} $<$ \ce{C6H5CH2+}
$\approx$ \ce{(CH3)3C+} $<$ \ce{CH3OCH2+}}: مجموعات الألكيل تعطي إلكترونات؛
والأليل والبنزيل يُلغيان تمركز الشحنة (بنيتان وأربع بنى)؛ \omterm{def:b1:lewis-resonance-vsepr:lewis}{والزوج الحر}
للأكسجين يعطي بنية فيها ثمانية كاملة على كل ذرة. وترتيب
وسط القائمة يتعلق بالوسط.
\end{solution}

\begin{solution}{exo:b1:electronic-effects:9}
يهاجم \ce{OH-} (\omterm{def:b1:electronic-effects:nucleophile}{محب للنواة}) كربون \ce{CH3Br} (\omterm{def:b1:electronic-effects:nucleophile}{محب للإلكترونات})،
ويغادر زوج C--Br على Br. ويعطي \ce{H2O} زوجًا للأيون \ce{H+} (وهنا الاسمان
المستعملان هما الحمض والقاعدة). ويهاجم \ce{CN-} كربون الكربونيل في
الإيثانال، وينتقل زوج $\pi$ للرابطة C=O إلى الأكسجين. ويعطي \ce{NH3} زوجه الحر
للمدار الفارغ للبور في \ce{BF3}.
\end{solution}

\begin{solution}{exo:b1:electronic-effects:10}
لا: $K_a = 10^{-0.51} = 0.31$ أكبر بكثير من $C$؛ فالحمض متفكك
تفككًا شبه تام. والمعادلة $x^2/(0.010 - x) = 0.31$ تعطي $x = \qty{9.7e-3}{mol/L}$،
$\mathrm{pH} = 2.01$، \omterm{def:b1:predominant-reaction:dissociation}{ودرجة التفكك} \qty{97}{\percent}.
\end{solution}

\begin{solution}{exo:b1:electronic-effects:11}
تدفع مجموعات الألكيل الإلكترونات ($+I$) نحو أكسجين الألكوكسيد فتجعل
شحنته أقل ثباتًا، وكلما كثرت المجموعات كان ذلك أشد. والألكوكسيدات الأكبر
أقل إذابةً أيضًا بالماء، الذي يثبّت ذرة أكسجين صغيرة مشحونة
أفضل من ذرة مدفونة بين مجموعات الألكيل. وكلا الأمرين يرفع
$\mathrm{p}K_a$.
\end{solution}

\begin{solution}{exo:b1:electronic-effects:12}
في الأميد يكون \omterm{def:b1:lewis-resonance-vsepr:lewis}{الزوج الحر} للنيتروجين مشتركًا مع C=O
(البنية \ce{N+=C-O^-})؛ وأخذه لبروتون أو \omterm{def:b1:electronic-effects:nucleophile}{لمحب للإلكترونات}
يهدم ذلك التثبيت. وفي \omterm{def:b1:predominant-reaction:strong-acid}{الحمض القوي} يتبرتن الأميد على
الأكسجين، الذي يحمل الشحنة الجزئية السالبة.
\end{solution}

\begin{solution}{pb:b1:electronic-effects:1}
\textbf{1.} $10^{1.89} = 78$؛ $10^{1.52} = 33$؛ $10^{0.84} = 6.9$.
\textbf{2.} لا: كل ذرة كلور إضافية تساعد أقل. فشحنة
الكربوكسيلات منتشرة جيدًا من قبل، وتصير القياسات أقل
دقة كلما اقترب الحمض من التفكك التام.
\textbf{3.} كل ذرة Cl تجذب الكثافة الإلكترونية عبر الروابط C--C وC--O
وتثبّت الكربوكسيلات ($-I$).
\textbf{4.} يتلاشى \omterm{def:b1:electronic-effects:inductive}{التأثير التحريضي} في حدود رابطتين أو ثلاث.
\textbf{5.} يبدو الحمضان متساويين في القوة (0.52 و0.51). ففي الماء،
كلاهما متفكك تفككًا شبه تام عند التراكيز المعتادة
(\cref{ch:b1:predominant-reaction}، التسوية): و$\mathrm{p}K_a$ القريبة من 0
هي حد ما يستطيع الماء تمييزه، والقيمتان تحملان
عدم يقين يبلغ بضعة أعشار.
\textbf{6.} حمض الإيثانويك: الصورة الحمضية (\qty{98}{\percent})؛ كلورو الإيثانويك:
الصورتان، مع غلبة طفيفة للأنيون ($10^{3 - 2.87} = 1.3$)؛ والثلاثة
الأخرى: الأنيون.
\textbf{7.} \ce{CH3-C(=O)-O^-} $\leftrightarrow$ \ce{CH3-C(-O^-)=O}.
\textbf{8.} رابطتان C--O متساويتان، وسط بين الرابطة الثنائية والأحادية.
\textbf{9.} الشحنة على ذرة الأكسجين الوحيدة، ولا رابطة $\pi$ تتقاسمها
معها.
\textbf{10.} $10^{15.9 - 4.76} = 10^{11.1}$، أي نحو $10^{11}$.
\textbf{11.} تنتشر شحنة الفينوكسيد إلى ثلاث ذرات كربون في الحلقة، أقل
كهرسلبيةً من الأكسجين: كسب، لكنه أصغر منه في الكربوكسيلات:
$\mathrm{p}K_a$ يساوي 9.99 بين 4.76 و15.9.
\textbf{12.} $K = 10^{6.37 - \mathrm{p}K_a}$: حمض الإيثانويك $10^{1.61} = 41$،
تفاعل مرجَّح (ينطلق ثنائي أكسيد الكربون مع فائض من الهيدروجين كربونات)؛
والفينول $10^{-3.62}$ والإيثانول $10^{-9.5}$: لا تفاعل.
\textbf{13.} \babelsublr{\foreignlanguage{arabic}{4-نيتروفينول (7.16)} $>$ \foreignlanguage{arabic}{2-نيتروفينول (7.23)} $>$ \foreignlanguage{arabic}{3-نيتروفينول
(8.36)} $>$ \foreignlanguage{arabic}{الفينول (9.99)}}.
\textbf{14.} البنية الكينونية في \cref{exo:b1:electronic-effects:3}،
والشحنة على أكسجين من مجموعة النيترو.
\textbf{15.} لا تصل الشحنة إلا إلى الكربونين \textit{ortho}
و\textit{para} بالنسبة إلى الأكسجين؛ وكربون النيترو في الموضع \textit{meta}.
\textbf{16.} تأثير $-I$ لمجموعة \ce{NO2}.
\textbf{17.} في 2-نيتروفينول تكوّن مجموعة OH \omterm{def:b1:intermolecular-forces-solvents:hbond}{رابطة هيدروجينية} مع أكسجين
من مجموعة النيترو المجاورة، وهذا يثبّت الصورة الحمضية ويجعل
نزع البروتون أصعب قليلًا.
\textbf{18.} كسور الأنيون $1/(1 + 10^{\mathrm{p}K_a - 7.5})$: 4-نيترو
\qty{69}{\percent}، و2-نيترو \qty{65}{\percent}، و3-نيترو
\qty{12}{\percent}. يمتص حمض 4-نيتروفينول وأنيونه امتصاصًا مختلفًا،
والأنيون في المرئي: فيظهر لونه حول قيمة $\mathrm{p}K_a$ له.
\textbf{19.} $x^2/(0.010 - x) = 10^{-0.52} = 0.30$: $x = \qty{9.7e-3}{mol/L}$،
$\mathrm{pH} = 2.01$: يكاد يكون \omterm{def:b1:predominant-reaction:strong-acid}{حمضًا قويًا}.
\textbf{20.} $\frac12(4.76 + 2.00) = 3.38$.
\textbf{21.} \ce{CF3COOH + CH3COO- -> CF3COO- + CH3COOH}، $K = 10^{4.76 -
0.52} = 10^{4.24}$: \omterm{def:b1:extent-q-and-k:quantitative}{تفاعل كمي}.
\textbf{22.} $K_a(\ce{CF3COOH})/K_a(\ce{CH3COOH}) =$ \textbf{$10^{4.24} =
\num{1.7e4}$}، أي نحو عشرة آلاف.
\end{solution}
